\documentclass[10pt]{article}

\usepackage[utf8]{inputenc}

\usepackage[T1]{fontenc}

\usepackage{amsmath}

\usepackage{amsfonts}

\usepackage{amssymb}

\usepackage[version=4]{mhchem}

\usepackage{stmaryrd}

\usepackage{mathrsfs}

\usepackage{bbold}

\usepackage{graphicx}

\usepackage[export]{adjustbox}

\graphicspath{ {./images/} }



\DeclareUnicodeCharacter{0131}{$\imath$}



\begin{document}

извольной размерности) классич. Римана-Роха теоремы (см. [2]). После создания выстей алгебраич. $K$-теории, то есть когомологич. теории функторов $K_{i}$, $i>0$ (см. [1], [12]) начинается интенсивное проникновение ее идей в алгебраич. геометрию. If настоящему времени можно выделить следующие области исследований в этом направлении.



\begin{enumerate}

  \item Изучение алгебраич. циклов на алгебраич. многообразиях. Пусть $X$ - гладкое алгебраич. многообразие, $C H(X)$ - кольцо Чжоу алгебраич. циклов на $X$ по модулю рациональной әквивалентности. Тогда имеют место изоморфизмы

\end{enumerate}



\[

\begin{gathered}

C H(X) \otimes \mathbb{Q} \cong K_{0}(X) \otimes \mathbb{Q}, \\

C H(X) \cong \bigcup_{i \geqslant 0} H^{i}\left(X, \mathscr{K}_{i}\right),

\end{gathered}

\]



где $\mathscr{K}_{i}$ - пучок в топологии Зариского, связанный с предпучком $U \rightarrow K_{i} \cdot \Gamma\left(U, \sigma_{X}\right)$. Эти факты являются освовой для изучения колец $C H(X)$ методами $K$-теории. В тастности, этими методами доказаны теоремы о конечности для групп Чжоу 0-циклов на арифмметич. поверхностях [4].



\begin{enumerate}

  \setcounter{enumi}{1}

  \item Значения дзета-функции и $L$-функции алгебраич. многообразий в целых точках. Имеются гипотезы о связи значений дзета-функций полей алгебраич. чисел в целых точках с порядками подгрупп кручения в $K$ Ф. их колец целых, а также значений $L$-функций многообразий над полями алгебраич. чисел в целых точках с рангами их групп $K_{i}$ и объемами решеток, порождаемых образом $K-\Phi$. в их когомологиях (cм. [1], [9]). Эти гипотезы подтверждаются в ряде частных случаев, они дополняют гипотезы Берча - Суиннертона-Дайера (см. Дзета-функция в алгебраической геометрии).



  \item Теория полей классов в высших размерностях описывает группу Галуа максимального абелева распирения полей рациональных функций арифметич. схем размерности $i \geqslant 1$, а также соответствующих локальных объектов (i-мерных локальных полей [7], [10]). В этом описании роль, к-рую обычно играет мультипликативная группа размерности 1 , выполняют групшы $K_{i}$ Милнора.



  \item Связь кристаллич. когомологий и деформаций K-Ф. (см. [3]).



  \item Теория характеристич. классов в алгебраич. $K$-теории и теорема Римана - Роха - Гротендика (cM. [5], [11]).



  \item Вычисление алгебраического $K$-Ф. для широкого класса схем. В частности, вычислены $K$-Ф. с конечными коәффициентами алгебраически замкнутых полей [8].



\end{enumerate}



Jum.: [1] Algebraic $K$-theory, v. 1, p. 85-147, v. 2, p. 489501, B.- - Hdlb. - N. Y., 1973; '[2] Théorie des intersections et théorème de Riemann-Roch, B.- Hdlb. - N. Y., 1971; [3]



\includegraphics[max width=\textwidth, center]{2023_07_09_73812cd618761a3c6749g-1}

"C. r. Acad. sci.", 1982 , t. 294 , ser. 1, p $749-52$; [5] G i l" e t t H. "Adv. in math.", 1981 , v 40, p. 203-89; [6] H a rr i s B., "S eg a l G., "Ann. math.", 1975, v. 101, № 1, p. 20-33; [7] K a t o K., «J. Fac. Sci. Univ. Tokyo", $1979-80$, Sec. IA,

v. 26 , № 2, p. $303-76$, v. 27 , № 3, p. $603-83$;



\includegraphics[max width=\textwidth, center]{2023_07_09_73812cd618761a3c6749g-1(1)}

[9] Б е й л и н с о н А. А., "Функциональный анализ", 1980 , т. 14 , в. 2, с. $46-47$, [10] П' а р Ш и н А. Н., «Докл. АН СССР», 1978 , т. 243 , № 4, c. 855-58, [11] ШI е х т м а Н B. В., "Успехи матем. наук», 1980 , т. 35 , в. 6 , с. $179-80 ;$ [12] Итоги науки и техники. Алгебра. Топология. Геометрия, т. 20 , М., 1982,

с. $71-152$.

В. В. Пехтман.



ФУНКТОРНЫЙ МОРФИЗМ - аналог понятия гомоморфизма (левых) модулей с общим кольцом скаляров (роль кольца при әтом играет область определения функторов, а сами функторы играют роль модулеӥ). Пусть $F_{1}$ и $F_{2}$ - одноместные ковариантные функторы из категории $\mathfrak{R}$ в категорию (5. Ф у н к т о р н ы м м о р ф з м о м $\varphi: F_{1} \longrightarrow F_{2}$ наз. такое сопоставление каждому объекту $A$ из $\mathfrak{A}$ морфизма $\varphi_{A}: F_{1}(A) \longrightarrow F_{2}(A)$, что для любого морфиима $\alpha: A \longrightarrow B$ из $\mathfrak{\Re}$ коммутативна следующая диаграмма:



\[

\begin{gathered}

F_{1}(A) \stackrel{F_{1}(\alpha)}{\longrightarrow} F_{1}(B) \\

\varphi_{A \downarrow} \downarrow \varphi_{B} \\

F_{2}(A) \underset{F_{2}(\alpha)}{\longrightarrow} F_{2}(B) .

\end{gathered}

\]



Если $F_{1}=F_{2}$, то, полагая $\varphi_{A}=1_{F_{1}(A)}$, получают т. н. тождест в енвы й морфизм функтора $F_{1}$. Если $\varphi: F_{1} \longrightarrow F_{2}$ и $\psi: F_{2} \longrightarrow F_{3} \longrightarrow$ два Ф. м., то, полагая $(\varphi \psi)_{A}=\varphi_{A} \psi_{A}$, получают $\Phi_{\text {м. }} \varphi \psi: F_{\mathbf{i}} \longrightarrow F_{3}$, называемый произведением $\varphi$ и $\psi$. Композıция $\Phi$. м. ассоциативна. Поэтому для малой категории $\mathfrak{R}$ все функторы из $\mathfrak{R}$ в $\sqrt{5}$ и их Ф. м. образуют т. н. к а т егорию фу ниторов Funct $(\Re,(5)$, или к а те гори и ди а грамм с о схемой $\mathfrak{\Re}$.



Пусть $\varphi: F_{1} \longrightarrow F_{2}: \mathfrak{R} \rightarrow(5-\Phi$. м. и $G: \mathfrak{M} \rightarrow \mathfrak{M}$, $H: \mathbb{S} \longrightarrow \mathfrak{N}$ - два функтора. Формулы



\[

\begin{aligned}

& \forall B \in \mathrm{Ob} \mathfrak{M}(G * \varphi)_{B}=\varphi_{G(B)}, \\

& \forall A \in \mathrm{Ob} \mathfrak{M}(\varphi * H)_{A}=H\left(\varphi_{A}\right)

\end{aligned}

\]



определяют Ф. м. $G * \varphi: G F_{1} \longrightarrow G F_{2}$ и $\varphi * H: F_{1} H \longrightarrow F_{2} H$ соответственно. Тогда для любых Ф. м. $\varphi: F_{1} \rightarrow F_{2}$ : $\mathfrak{A} \longrightarrow\left(5\right.$ и $\psi: H_{1} \longrightarrow H_{2}:(5) \Re$ справедливо соотнотение



\[

\left(\varphi * H_{1}\right)\left(F_{2} * \psi\right)=\left(F_{1} * \psi\right)\left(\varphi * H_{2}\right)

\]



Ф.м. наз. также е с т е с т в е н н ы м и п р е о б р азов ани я м и ф ун т о ров. По аналогии с Ф. м. одноместных функторов определяются Ф. м. многоместных функторов. $М$. ШІ. Цаленко.



ФУНКЦИЙ ДЕИСТВИТЕЛЬНОГО ПЕРЕМЕННОГО ТЕОРИЯ - область математич. анализа, в к-рой изучаются вопросы представления и приближения функций, их локальные и глобальные свойства. Для современной Ф. д. п. т. характерно широкое применение теоретико-множественных методов наряду, естественно, с классическими.



Таким образом, объектом изучения в Ф. Д. П. т. является функция. По поводу әтого понятия Н. Н. Лузин [3] писал: «Оно не сложилось сразу, но, возникнув более двухсот лет тому назад в знаменитом споре о звучащей струне, подверглось глубоким изменениям уже в начавшейся тогда энергичной полемике. С тех пор идут непрестанное углубление и эволюция этого понятия, которые продолжаются до настоящего времени. Поэтому ни одно отдельное формальное определение не может охватить всего содержания этого понятия...». В соответствии с этим представляется вполне естественным отнести истоки зарождения Ф. Д. П. т. ко времени спора о колеблющейся струне [Л. Эйлер (L. Euler), Д. Бернулли (D. Bernoulli), Ж. Д'Аламбер (J. D'Alembert), Ж. Лагранж (J. Lagrange) и др.], хотя становление әтой теории происходило в 19 в. [Ж. Фурье (J. Fourier), О. Коши (А. Cauchy), Н. И. Лобачевский, П. Дирихле (Р. Dirichlet), Б. Риман (B. Riemann), П. Л. Чебышев, К. Жордан (C. Jordan) и др.].



В классич. анализе изучались в основном функции, имеющие определенную степень гладкости. Однако во 2-й пол. 19 в. в анализе четко обрисовались нек-рые проблемы, ждавшие своего решения и касающиеся более общих классов функций, а также более глубокого изучения и гладких функций. $\mathfrak{K}$ таким проблемам следует отнести проблемы меры множества, длин кривых и площалей поверхностей, приближения и представления функций, первообразной и интеграла, взаимосвязи интегрирования и дифференцирования, почленное интегрирование и дифференцирование рядов, свойства функций, полученных в результате предельного перехода и др. Решение этих проблем имело принципиальное значение для математики. Классич. методы анализа





\end{document}